\documentclass[aspectratio=169,11pt]{beamer}

% ==============================
% NechCode - Analisa Ide Kelayakan
% Wedding Invitation + Gifting Site berbasis AI
% ==============================

\usepackage[T1]{fontenc}
\usepackage[utf8]{inputenc}
\usepackage[indonesian]{babel}
\usepackage{lmodern}
\usepackage{amsmath,amssymb}
\usepackage{booktabs,tabularx,array,multirow}
\usepackage{tikz}
\usepackage{pifont}
\usepackage[table]{xcolor}
\usepackage{ragged2e}
\usepackage{graphicx}

\usetikzlibrary{positioning,calc,shadows.blur,arrows.meta}

% ---------- Colors ----------
\definecolor{Navy}{HTML}{143D73}
\definecolor{Blue}{HTML}{2F6ECB}
\definecolor{SoftBlue}{HTML}{EAF2FF}
\definecolor{Cream}{HTML}{FAF4EA}
\definecolor{Sand}{HTML}{E8D8C2}
\definecolor{TextDark}{HTML}{243447}
\definecolor{TextMuted}{HTML}{6B7A90}
\definecolor{Good}{HTML}{57B65F}
\definecolor{Bad}{HTML}{E76666}
\definecolor{Gold}{HTML}{D3B16A}
\definecolor{LightGreen}{HTML}{DFF7E7}
\definecolor{LightYellow}{HTML}{F7F1CC}
\definecolor{LightBrown}{HTML}{E9DDC7}

% ---------- Global Style ----------
\setbeamersize{text margin left=0.65cm,text margin right=0.65cm}
\setbeamertemplate{navigation symbols}{}
\setbeamertemplate{itemize item}{\color{Blue}$\blacktriangleright$}
\setbeamertemplate{itemize subitem}{\color{Blue}$\bullet$}
\setbeamercolor{normal text}{fg=TextDark,bg=white}
\setbeamercolor{structure}{fg=Navy}
\setbeamercolor{frametitle}{fg=Navy,bg=Cream}
\setbeamercolor{title}{fg=Navy}
\setbeamercolor{subtitle}{fg=Blue}
\setbeamercolor{block title}{fg=white,bg=Navy}
\setbeamercolor{block body}{fg=TextDark,bg=SoftBlue}
\setbeamerfont{title}{size=\LARGE,series=\bfseries}
\setbeamerfont{subtitle}{size=\large,series=\bfseries}
\setbeamerfont{frametitle}{size=\Large,series=\bfseries}
\setbeamerfont{block title}{series=\bfseries}

% ---------- Frame title ----------
\makeatletter
\setbeamertemplate{frametitle}{%
  \vspace{0.12cm}
  \begin{beamercolorbox}[wd=\paperwidth,ht=1cm,dp=0.2cm,leftskip=0.7cm,rightskip=0.25cm]{frametitle}
    {\usebeamerfont{frametitle}\insertframetitle}
    \ifx\insertframesubtitle\@empty
    \else
      \\[-1pt]{\small\color{TextMuted}\insertframesubtitle}
    \fi
  \end{beamercolorbox}
}
\makeatother

% ---------- Footer ----------
\setbeamertemplate{footline}{%
  \leavevmode
  \hbox{%
    \begin{beamercolorbox}[wd=.75\paperwidth,ht=3ex,dp=1.2ex,leftskip=0.55cm]{author in head/foot}%
      \color{TextMuted}\scriptsize NechCode -- Ide Bisnis dan Analisa Kelayakan
    \end{beamercolorbox}%
    \begin{beamercolorbox}[wd=.25\paperwidth,ht=3ex,dp=1.2ex,rightskip=0.55cm]{date in head/foot}%
      \hfill\color{TextMuted}\scriptsize \insertframenumber/\inserttotalframenumber
    \end{beamercolorbox}%
  }
}

% ---------- Helpers ----------
\newcommand{\pill}[1]{%
  \tikz[baseline=(X.base)] \node (X) [rounded corners=8pt, fill=SoftBlue, text=Navy, inner xsep=8pt, inner ysep=4pt, font=\scriptsize\bfseries] {#1};%
}

\newcommand{\yesicon}{%
  \tikz[baseline=-0.5ex]\node[circle,fill=Good,inner sep=1.2pt]{\color{white}\scriptsize\ding{51}};
}
\newcommand{\noicon}{%
  \tikz[baseline=-0.5ex]\node[circle,fill=Bad,inner sep=1.3pt]{\color{white}\scriptsize\ding{55}};
}
\newcommand{\midicon}{%
  \tikz[baseline=-0.5ex]\node[circle,fill=Gold,inner sep=1.3pt]{\color{white}\scriptsize\textbf{!}};
}

\newcommand{\statcard}[3]{%
\begin{tikzpicture}
\node[rounded corners=12pt, fill=Cream, draw=Sand, line width=0.7pt,
minimum width=3.1cm, minimum height=1.8cm, align=left, inner sep=8pt] {
\begin{minipage}{2.45cm}
{\color{TextMuted}\scriptsize #1}\\[1pt]
{\color{Navy}\Large\bfseries #2}\\[2pt]
{\color{TextDark}\scriptsize #3}
\end{minipage}
};
\end{tikzpicture}%
}

\newcommand{\featurecard}[2]{%
\begin{tikzpicture}
\node[rounded corners=12pt, fill=white, draw=SoftBlue, line width=1pt,
text width=0.92\linewidth, inner sep=10pt, align=left] {
{\color{Navy}\bfseries #1}\\[2pt]
{\color{TextDark}\small #2}
};
\end{tikzpicture}%
}

\newcommand{\riskcard}[3]{%
\begin{tikzpicture}
\node[rounded corners=12pt, fill=#1!10, draw=#1!45, line width=0.8pt,
text width=0.92\linewidth, inner sep=9pt, align=left] {
{\color{Navy}\bfseries #2}\\[2pt]
{\color{TextDark}\small #3}
};
\end{tikzpicture}%
}

\newcommand{\outlinebox}[3]{%
\begin{tikzpicture}
\node[rounded corners=4pt, fill=#1, draw=black!35, line width=0.7pt,
minimum width=5.35cm, minimum height=1.9cm, align=center] {
{\color{black}\bfseries\itshape #2}\\[1pt]
{\color{black}\bfseries\itshape #3}
};
\end{tikzpicture}%
}

% ---------- Metadata ----------
\title{Ide Bisnis dan Analisa Kelayakan: \\ Undangan Nikah Digital dan Gifting Site berbasis AI}
\subtitle{Fokus khusus: analisa ide kelayakan NechCode}
\author{Kelompok 12}
\date{Maret 2026}

\begin{document}

% =====================================
% Title Slide
% =====================================
\begin{frame}[plain]
\begin{tikzpicture}[remember picture,overlay]
  %% Full Navy background
  \fill[Navy] (current page.south west) rectangle (current page.north east);

  %% Large atmospheric circles - right decorative zone
  \fill[Blue!50] ([xshift=14.3cm,yshift=8.5cm]current page.south west) circle (2.6cm);
  \fill[Blue!28] ([xshift=15.9cm,yshift=2.8cm]current page.south west) circle (2.4cm);
  \fill[Blue!17] ([xshift=13.5cm,yshift=5.0cm]current page.south west) circle (4.0cm);

  %% Gold accent dots (premium touch)
  \fill[Gold] ([xshift=11.5cm,yshift=7.9cm]current page.south west) circle (0.18cm);
  \fill[Gold!72] ([xshift=12.2cm,yshift=7.4cm]current page.south west) circle (0.12cm);
  \fill[Gold!50] ([xshift=10.8cm,yshift=8.2cm]current page.south west) circle (0.09cm);

  %% Subtle white dots
  \fill[white!16] ([xshift=13.3cm,yshift=6.7cm]current page.south west) circle (0.14cm);
  \fill[white!11] ([xshift=14.2cm,yshift=5.7cm]current page.south west) circle (0.09cm);

  %% Left accent strip
  \fill[Blue] (current page.south west)
               rectangle ([xshift=0.28cm]current page.north west);

  %% Bottom bar
  \fill[Blue!80] (current page.south west)
                  rectangle ([yshift=0.40cm]current page.south east);

  %% Floating feature badges (right zone, layered on circles)
  \node[rounded corners=10pt, fill=Navy!85, draw=SoftBlue!55, line width=0.7pt,
        text=white!88, font=\scriptsize\bfseries, inner xsep=10pt, inner ysep=5pt]
    at ([xshift=12.5cm,yshift=7.0cm]current page.south west) {AI};
  \node[rounded corners=10pt, fill=Navy!85, draw=SoftBlue!55, line width=0.7pt,
        text=white!88, font=\scriptsize\bfseries, inner xsep=10pt, inner ysep=5pt]
    at ([xshift=14.5cm,yshift=6.1cm]current page.south west) {RSVP};
  \node[rounded corners=10pt, fill=Navy!85, draw=SoftBlue!55, line width=0.7pt,
        text=white!88, font=\scriptsize\bfseries, inner xsep=10pt, inner ysep=5pt]
    at ([xshift=13.3cm,yshift=3.0cm]current page.south west) {Gifting};

  %% Vertical NechCode label (far right)
  \node[anchor=center, text=white!22, font=\small\bfseries, rotate=90]
    at ([xshift=15.5cm,yshift=4.5cm]current page.south west) {NechCode};
\end{tikzpicture}

\vspace{1.35cm}
\hspace*{0.7cm}
\begin{minipage}[t]{9.2cm}
{\color{SoftBlue}\scriptsize\bfseries PROPOSAL BISNIS DIGITAL}\\[0.50cm]
{\color{white}\fontsize{17}{22}\selectfont\bfseries\raggedright
  Undangan Nikah Digital\\
  \& Gifting Site Berbasis AI}\\[0.10cm]
{\color{SoftBlue!80}\normalsize\raggedright untuk NechCode}\\[0.24cm]
{\color{Gold}\rule{5.0cm}{0.45pt}}\\[0.20cm]
{\color{white!72}\small\raggedright
  Analisa kelayakan: potensi pasar, diferensiasi produk,\\
  dan langkah MVP yang realistis.}\\[0.34cm]
\tikz[baseline=(X.base)]\node(X)[rounded corners=7pt,fill=Blue!65,text=white,inner xsep=8pt,inner ysep=4pt,font=\scriptsize\bfseries]{Wedding Tech};\hspace{0.06cm}%
\tikz[baseline=(X.base)]\node(X)[rounded corners=7pt,fill=Blue!65,text=white,inner xsep=8pt,inner ysep=4pt,font=\scriptsize\bfseries]{AI Template};\hspace{0.06cm}%
\tikz[baseline=(X.base)]\node(X)[rounded corners=7pt,fill=Blue!65,text=white,inner xsep=8pt,inner ysep=4pt,font=\scriptsize\bfseries]{RSVP};\hspace{0.06cm}%
\tikz[baseline=(X.base)]\node(X)[rounded corners=7pt,fill=Blue!65,text=white,inner xsep=8pt,inner ysep=4pt,font=\scriptsize\bfseries]{Gifting};\\[0.36cm]
{\color{white}\bfseries Kelompok 12}\hspace{0.25cm}{\color{SoftBlue!75}\small$\cdot$\hspace{0.08cm}Maret 2026}
\end{minipage}
\end{frame}

% =====================================
% Ringkasan bisnis
% =====================================
\begin{frame}[shrink=12]{Ringkasan Bisnis}{Konsep singkat NechCode}
\begin{columns}[T,totalwidth=\textwidth]
\column{0.56\textwidth}
\small
\begin{block}{One-line concept}
NechCode adalah layanan digital yang membantu pengguna membuat undangan nikah digital dan gifting site berbasis AI-template, lalu memantau RSVP dan status hadir tamu dengan lebih praktis.
\end{block}

\vspace{0.08cm}
\footnotesize
\begin{itemize}
  \item \textbf{Fokus layanan utama:} undangan nikah digital
  \item \textbf{Layanan pendukung:} gifting site untuk ulang tahun, wisuda, anniversary, dan appreciation
  \item \textbf{Target awal:} pasangan muda dan pengguna personal digital-native
  \item \textbf{Model singkat:} paket sekali beli, template premium, add-on RSVP dan gifting
\end{itemize}

\column{0.40\textwidth}
\featurecard{Nilai utama}{Cepat dibuat, lebih personal, mudah dibagikan, dan lebih rapi untuk memantau RSVP dibanding undangan manual.}\\[0.10cm]
\featurecard{Arah MVP}{User cukup isi nama, kontak, detail acara, pilih template, lalu sistem generate undangan personal dan dashboard tamu sederhana.}
\end{columns}
\end{frame}

% =====================================
% Bentuk analisa kelayakan
% =====================================
\begin{frame}{Bentuk Analisa Kelayakan}
\vspace{0.25cm}
\begin{center}
\begin{tikzpicture}
\node at (0,0) {};
\node[anchor=north west] (a) at (-5.7,1.7) {\outlinebox{LightYellow}{Product/Service Feasibility}{}};
\node[anchor=north west] (b) at (0.2,1.7) {\outlinebox{LightYellow}{Industry/Target Market}{Feasibility}};
\node[anchor=north west] (c) at (-5.7,-0.6) {\outlinebox{LightYellow}{Organizational Feasibility}{}};
\node[anchor=north west] (d) at (0.2,-0.6) {\outlinebox{LightYellow}{Financial Feasibility}{}};
\end{tikzpicture}
\end{center}

\vspace{0.3cm}
\begin{center}
\pill{Sesuai outline kelas}\hspace{0.15cm}
\pill{Dipakai sebagai kerangka analisa deck ini}
\end{center}
\end{frame}

% =====================================
% Outline komprehensif
% =====================================
\begin{frame}{Outline Analisa Kelayakan (1/2)}{Kerangka yang digunakan}
\small
\renewcommand{\arraystretch}{1.28}
\setlength{\tabcolsep}{5pt}
\begin{center}
\begin{tabularx}{\textwidth}{|>{\RaggedRight\arraybackslash}p{5.2cm}|X|}
\hline
\rowcolor{yellow!70}
\textbf{Product/Service Feasibility\newline (Kelayakan Produk/Jasa)} &
A. Keinginan atas produk/jasa\\
& B. Permintaan atas produk/jasa\\
\hline
\rowcolor{blue!10}
\textbf{Industry/Target Market Feasibility\newline (Kelayakan Industri/Target Pasar)} &
A. Kemenarikan industri\\
& B. Kemenarikan target pasar\\
\hline
\rowcolor{green!25}
\textbf{Organizational Feasibility\newline (Kelayakan Organisasi)} &
A. Kecakapan organisasi\\
& B. Kecukupan sumberdaya\\
\hline
\end{tabularx}
\end{center}
\end{frame}

\begin{frame}{Outline Analisa Kelayakan (2/2)}{Kerangka yang digunakan}
\small
\renewcommand{\arraystretch}{1.28}
\setlength{\tabcolsep}{5pt}
\begin{center}
\begin{tabularx}{\textwidth}{|>{\RaggedRight\arraybackslash}p{5.2cm}|X|}
\hline
\rowcolor{LightYellow}
\textbf{Financial Feasibility\newline (Kelayakan Finansial)} &
A. Sejumlah uang yang dibutuhkan untuk memulai bisnis\\
& B. Performa finansial bisnis yang sejenis\\
& C. Keseluruhan kemenarikan finansial dari bisnis yang diusulkan\\
\hline
\rowcolor{LightBrown}
\textbf{Overall Assessment\newline (Penilaian Keseluruhan)} & Ringkasan penilaian akhir ide bisnis\\
\hline
\end{tabularx}
\end{center}

\vspace{0.35cm}
\begin{center}
\pill{Sesuai outline kelas}\hspace{0.15cm}
\pill{Dipakai sebagai kerangka analisa deck ini}
\end{center}
\end{frame}

% =====================================
% Observing Trend
% =====================================
\begin{frame}{Observing Trend}{Mengapa ide ini layak diamati}
\begin{columns}[T,totalwidth=\textwidth]
\column{0.59\textwidth}
\small
\begin{itemize}
  \item Pengguna digital Indonesia semakin terbiasa menerima link acara, RSVP online, dan berbagi momen lewat WhatsApp atau media sosial.
  \item Pasangan muda cenderung mencari solusi undangan yang praktis, cepat, dan hemat dibanding undangan cetak penuh.
  \item AI makin relevan untuk mempercepat penulisan copy, personalisasi, dan rekomendasi desain.
  \item Gifting digital juga makin masuk akal karena orang ingin memberi pengalaman spesial yang bisa dibagikan secara online.
\end{itemize}

\column{0.39\textwidth}
\begin{tikzpicture}
  \node[rounded corners=10pt, fill=Cream, draw=Sand, line width=0.7pt,
    text width=2.05cm, inner sep=7pt, align=left] (sa) at (0,0) {
    {\color{TextMuted}\scriptsize Digital Sharing}\\[1pt]
    {\color{Navy}\large\bfseries Kuat}\\[2pt]
    {\color{TextDark}\scriptsize Cocok untuk produk berbasis link.}
  };
  \node[rounded corners=10pt, fill=Cream, draw=Sand, line width=0.7pt,
    text width=2.05cm, inner sep=7pt, align=left, right=0.18cm of sa] {
    {\color{TextMuted}\scriptsize Mobile First}\\[1pt]
    {\color{Navy}\large\bfseries Tinggi}\\[2pt]
    {\color{TextDark}\scriptsize Undangan mudah dikirim lewat WhatsApp.}
  };
\end{tikzpicture}

\vspace{0.16cm}
\featurecard{Kenapa menarik?}{Saat personalisasi, distribusi digital, dan AI bertemu dalam satu produk, diferensiasi NechCode menjadi lebih jelas.}
\end{columns}

\vspace{0.14cm}
\featurecard{Makna untuk NechCode}{Produk bukan hanya menjual template undangan, tetapi pengalaman digital untuk momen spesial yang lebih praktis, personal, dan mudah dipantau.}
\end{frame}

% =====================================
% Solving problem
% =====================================
\begin{frame}{Solving Problem (1/2)}{Pain point yang ingin diselesaikan}
\begin{columns}[T,totalwidth=\textwidth]
\column{0.49\textwidth}
\begin{block}{Pain point undangan nikah digital}
\small
\begin{itemize}
  \item Sulit membuat undangan yang cepat tapi tetap estetik.
  \item Data tamu dan RSVP masih sering diatur manual.
  \item Personalisasi nama tamu biasanya butuh kerja tambahan.
  \item Banyak solusi terasa generik atau terlalu rumit.
\end{itemize}
\end{block}

\column{0.49\textwidth}
\begin{block}{Pain point gifting site}
\small
\begin{itemize}
  \item Pengguna ingin halaman hadiah yang spesial, bukan sekadar chat biasa.
  \item Tidak semua orang bisa desain mini-site sendiri.
  \item Sulit menggabungkan ucapan, gift request, dan halaman acara dalam satu alur.
  \item Pengalaman gifting sering belum terintegrasi dengan momen event.
\end{itemize}
\end{block}
\end{columns}
\end{frame}

\begin{frame}{Solving Problem (2/2)}{Nilai yang ditawarkan NechCode}
\vspace{0.4cm}
\featurecard{Nilai praktis untuk pelanggan}{NechCode membantu user membuat undangan atau gifting page yang cepat jadi, lebih personal, dan lebih mudah dipantau daripada proses manual.}

\vspace{0.3cm}
\begin{columns}[T,totalwidth=\textwidth]
\column{0.33\textwidth}
\begin{tikzpicture}
\node[rounded corners=12pt, fill=Good!12, draw=Good!40, line width=0.8pt,
text width=0.92\linewidth, inner sep=10pt, align=left] {
{\color{Navy}\bfseries Undangan}\\[3pt]
{\color{TextDark}\small Cepat dibuat, estetik, tamu bisa RSVP langsung dari link.}
};
\end{tikzpicture}

\column{0.33\textwidth}
\begin{tikzpicture}
\node[rounded corners=12pt, fill=Blue!8, draw=Blue!30, line width=0.8pt,
text width=0.92\linewidth, inner sep=10pt, align=left] {
{\color{Navy}\bfseries Gifting}\\[3pt]
{\color{TextDark}\small Halaman gift yang spesial, mudah dibagikan dan dikelola.}
};
\end{tikzpicture}

\column{0.33\textwidth}
\begin{tikzpicture}
\node[rounded corners=12pt, fill=Gold!12, draw=Gold!40, line width=0.8pt,
text width=0.92\linewidth, inner sep=10pt, align=left] {
{\color{Navy}\bfseries AI Assist}\\[3pt]
{\color{TextDark}\small Copy, wording, dan desain dibantu AI agar lebih personal.}
};
\end{tikzpicture}
\end{columns}
\end{frame}

% =====================================
% Gap pasar
% =====================================
\begin{frame}{Gap Pasar}{Celah yang ingin dimasuki NechCode}
\begin{columns}[T,totalwidth=\textwidth]
\column{0.54\textwidth}
\begin{itemize}
  \item Template umum memang banyak, tetapi user masih harus merancang isi, style, dan alur tamu sendiri.
  \item Jasa custom lebih personal, namun lebih mahal dan lebih lama.
  \item Platform undangan existing kuat di wedding, tetapi belum semuanya punya integrasi gifting yang menarik.
  \item Ada ruang untuk solusi yang lebih ringan daripada jasa custom, namun lebih terarah daripada tool desain umum.
\end{itemize}

\column{0.40\textwidth}
\begin{tikzpicture}
  \node[rounded corners=12pt, fill=Navy, inner sep=12pt, text width=0.92\linewidth] {
    \color{white}\textbf{Gap terkuat}\\[4pt]
    \color{SoftBlue}\small Kombinasi template AI, personalisasi tamu, RSVP tracking, dan gifting page dalam satu alur sederhana.
  };
\end{tikzpicture}
\end{columns}
\end{frame}

% =====================================
% Competitor matrix
% =====================================
\begin{frame}{Analisa Kompetitor}{Matriks perbandingan fitur}
\footnotesize
\renewcommand{\arraystretch}{1.05}
\setlength{\tabcolsep}{5.2pt}
\begin{center}
\begin{tabularx}{\textwidth}{>{\RaggedRight\arraybackslash}X >{\centering\arraybackslash}m{1.75cm} >{\centering\arraybackslash}m{1.6cm} >{\centering\arraybackslash}m{1.6cm} >{\centering\arraybackslash}m{1.75cm}}
\rowcolor{Cream}
\textbf{Perbandingan Fitur} & \textbf{Wevitation} & \textbf{Wedew} & \textbf{Invitanku} & \textbf{NechCode} \\
\midrule
Template undangan digital & \yesicon & \yesicon & \yesicon & \yesicon \\
RSVP tracking & \yesicon & \yesicon & \yesicon & \yesicon \\
Share mudah ke WhatsApp & \yesicon & \yesicon & \yesicon & \yesicon \\
AI bantu copy / desain & \noicon & \noicon & \midicon & \yesicon \\
Integrasi gifting site & \midicon & \noicon & \noicon & \yesicon \\
Personalisasi bahasa lokal & \midicon & \midicon & \yesicon & \yesicon \\
Dashboard sederhana dan ringan & \yesicon & \midicon & \yesicon & \yesicon \\
Potensi bundling wedding + hadiah & \midicon & \noicon & \noicon & \yesicon \\
Diferensiasi produk yang jelas & \midicon & \midicon & \midicon & \yesicon \\
\end{tabularx}
\end{center}

\vspace{0.15cm}
\begin{columns}[T,totalwidth=\textwidth]
\column{0.66\textwidth}
\featurecard{Insight utama}{Ruang diferensiasi NechCode paling kuat ada pada kombinasi \textbf{AI assistance}, \textbf{gifting integration}, dan pengalaman pengguna yang tetap ringan serta mudah dipahami.}

\column{0.32\textwidth}
{\scriptsize
\yesicon~Tersedia\\[2pt]
\midicon~Sebagian / belum kuat\\[2pt]
\noicon~Belum menjadi fokus
}
\end{columns}
\end{frame}

% =====================================
% Product/service feasibility
% =====================================
\begin{frame}{Product/Service Feasibility}{Kelayakan produk/jasa}
\begin{columns}[T,totalwidth=\textwidth]
\column{0.49\textwidth}
\begin{block}{A. Keinginan atas produk/jasa}
\begin{itemize}
  \item Produk ini menarik karena menyederhanakan pembuatan undangan digital.
  \item Fitur RSVP dan personalisasi tamu memberi manfaat nyata.
  \item Gifting page menambah nilai emosional untuk momen spesial.
  \item Cocok dengan perilaku user yang ingin praktis namun tetap estetik.
\end{itemize}
\end{block}

\column{0.49\textwidth}
\begin{block}{B. Permintaan atas produk/jasa}
\begin{itemize}
  \item Permintaan awal paling kuat ada pada use case wedding invitation.
  \item Gifting site lebih cocok sebagai cross-sell atau pengembangan tahap berikutnya.
  \item Willingness to pay akan dipengaruhi desain, kemudahan, dan tracking tamu.
  \item Produk ini lebih mudah diuji bila dimulai dari fitur inti dulu.
\end{itemize}
\end{block}
\end{columns}
\end{frame}

% =====================================
% Industry/Target market feasibility
% =====================================
\begin{frame}{Industry/Target Market Feasibility}{Kelayakan industri dan target pasar}
\begin{columns}[T,totalwidth=\textwidth]
\column{0.49\textwidth}
\begin{block}{A. Kemenarikan industri}
\begin{itemize}
  \item Industri layanan digital event cukup menarik karena biaya awal relatif ringan.
  \item Distribusi produk berbasis web/link memudahkan scale.
  \item Hambatan masuk tidak terlalu tinggi, sehingga diferensiasi jadi faktor penting.
  \item Kompetisi cukup padat, tetapi pasar sudah teredukasi.
\end{itemize}
\end{block}

\column{0.49\textwidth}
\begin{block}{B. Kemenarikan target pasar}
\begin{itemize}
  \item Pasangan muda adalah entry market paling realistis.
  \item Mereka punya kebutuhan yang jelas, event-driven, dan willingness to pay lebih kuat.
  \item Segmen gifting personal dapat menyusul setelah wedding traction awal.
  \item Pasar cukup spesifik untuk MVP dan tidak terlalu luas.
\end{itemize}
\end{block}
\end{columns}
\end{frame}

% =====================================
% Organizational feasibility
% =====================================
\begin{frame}{Organizational Feasibility}{Kelayakan organisasi}
\begin{columns}[T,totalwidth=\textwidth]
\column{0.49\textwidth}
\begin{block}{A. Kecakapan organisasi}
\begin{itemize}
  \item Bisnis ini cocok dijalankan tim kecil.
  \item Skill inti: web development, UI design, product thinking, support, dan digital marketing.
  \item MVP tidak membutuhkan organisasi besar jika ruang lingkup dijaga sempit.
  \item Kompleksitas utama ada pada menjaga pengalaman pengguna tetap sederhana.
\end{itemize}
\end{block}

\column{0.49\textwidth}
\begin{block}{B. Kecukupan sumberdaya}
\begin{itemize}
  \item Yang wajib ada sejak awal: web app dasar, template awal, RSVP dashboard, landing page.
  \item Yang bisa menyusul: custom domain, AI layer lebih dalam, analytics, dan gift integration yang lebih kompleks.
  \item Resource ini relatif realistis untuk startup tahap awal.
\end{itemize}
\end{block}
\end{columns}
\end{frame}

% =====================================
% Financial feasibility
% =====================================
\begin{frame}[shrink=10]{Financial Feasibility (1/2)}{A \& B -- Biaya dan perbandingan bisnis sejenis}
\begin{columns}[T,totalwidth=\textwidth]
\column{0.49\textwidth}
\begin{block}{A. Uang yang dibutuhkan untuk memulai}
\begin{itemize}
  \item Domain dan hosting
  \item Development MVP
  \item Desain template awal
  \item Tools AI dan operasional dasar
  \item Marketing awal dan support sederhana
\end{itemize}
\end{block}

\column{0.49\textwidth}
\begin{block}{B. Performa finansial bisnis sejenis}
\begin{itemize}
  \item Model umum: freemium, paket sekali beli, premium template, add-on.
  \item Produk digital punya potensi margin baik jika custom work dijaga rendah.
\end{itemize}
\end{block}
\end{columns}

\vspace{0.22cm}
\begin{center}
\statcard{Model Awal}{One-time}{Paket per acara}\hspace{0.4cm}
\statcard{Upsell}{Add-on}{Gift, gallery, premium theme}\hspace{0.4cm}
\statcard{Risiko}{Harga}{Pasar sensitif harga}
\end{center}
\end{frame}

\begin{frame}{Financial Feasibility (2/2)}{C -- Kemenarikan finansial keseluruhan}
\begin{block}{C. Keseluruhan kemenarikan finansial}
\begin{itemize}
  \item Menarik karena biaya awal tidak terlalu berat.
  \item Pendapatan paling realistis: paket acara sekali beli.
  \item Gifting site dapat memperbesar average order value.
  \item Risiko utama ada pada harga dan tekanan kompetisi.
\end{itemize}
\end{block}

\vspace{0.22cm}
\featurecard{Kesimpulan finansial}{Produk ini realistis dibangun dengan anggaran terbatas. Fokus pada paket undangan inti dulu, lalu optimalkan average order value lewat add-on gifting dan premium template.}
\end{frame}

% =====================================
% Overall Assessment
% =====================================
\begin{frame}[shrink=5]{Overall Assessment}{Penilaian keseluruhan ide bisnis}
\begin{columns}[T,totalwidth=\textwidth]
\column{0.33\textwidth}
\riskcard{Good}{Kekuatan utama}{Pasar sudah teredukasi, use case nyata, dan distribusi digital sangat cocok untuk produk berbasis undangan dan gifting.}\\[0.09cm]
\riskcard{Bad}{Ancaman terbesar}{Kompetitor existing, template gratis, dan switching cost pengguna yang rendah.}

\column{0.33\textwidth}
\riskcard{Gold}{Kelemahan utama}{Jika terlalu luas antara wedding, gifting, dan AI, positioning bisa kabur dan eksekusi tim kecil menjadi berat.}\\[0.09cm]
\riskcard{Gold}{Kesimpulan awal}{Ide ini \textbf{layak diuji lebih lanjut} bila MVP difokuskan dulu pada undangan nikah digital dengan RSVP tracking.}

\column{0.33\textwidth}
\riskcard{Good}{Peluang terbesar}{Menjadi solusi yang lebih simpel dari platform wedding kompleks dan lebih terarah dari tool desain umum.}\\[0.09cm]
\riskcard{Good}{Rekomendasi}{Bangun bertahap: mulai dari invitation core, lalu AI assistance, lalu gifting integration.}
\end{columns}
\end{frame}

% =====================================
% Go to market
% =====================================
\begin{frame}{Strategi Masuk Pasar}{Urutan eksekusi yang lebih aman}
\begin{columns}[T,totalwidth=\textwidth]
\column{0.33\textwidth}
\featurecard{Tahap 1 -- MVP}{Undangan digital, RSVP, galeri foto, dan share link yang mudah. Fokus pada pengalaman dasar yang stabil dan estetik.}

\column{0.33\textwidth}
\featurecard{Tahap 2 -- AI Layer}{Tambahkan generator copy, rekomendasi wording, dan bantuan personalisasi desain agar produk lebih menonjol.}

\column{0.33\textwidth}
\featurecard{Tahap 3 -- Gifting}{Masukkan gift catalog atau gift request sederhana untuk memperluas monetisasi dan diferensiasi.}
\end{columns}

\vspace{0.28cm}
\begin{block}{Kenapa urutan ini penting?}
Dengan pendekatan bertahap, tim bisa memvalidasi minat pasar lebih dulu tanpa langsung terbebani kompleksitas operasional gifting.
\end{block}
\end{frame}

% =====================================
% Conclusion
% =====================================
\begin{frame}{Kesimpulan Analisa Kelayakan}{Ringkas, jelas, dan mudah dipahami audiens}
\begin{columns}[T,totalwidth=\textwidth]
\column{0.58\textwidth}
\begin{itemize}
  \item Ide ini \textbf{layak dikembangkan} karena kebutuhan pasar nyata, distribusi digital kuat, dan ada ruang diferensiasi.
  \item Keunikan NechCode ada pada \textbf{gabungan invitation platform + AI assistance + gifting experience}.
  \item Strategi paling aman adalah membuat MVP yang ringan dulu, lalu menambah AI dan gifting secara bertahap.
  \item Jika positioning dipertajam, produk ini berpotensi lebih menonjol dibanding layanan undangan digital yang hanya bermain di template.
\end{itemize}

\column{0.38\textwidth}
\begin{tikzpicture}
\node[rounded corners=16pt, fill=Navy, inner sep=12pt, text width=0.92\linewidth] {
\color{white}\bfseries Verdict\\[4pt]
\color{SoftBlue}\small \textbf{Layak untuk diuji sebagai MVP}, terutama bila NechCode fokus pada kemudahan penggunaan, personalisasi, dan diferensiasi gifting berbasis AI.
};
\end{tikzpicture}
\end{columns}
\end{frame}

% =====================================
% Team
% =====================================
\begin{frame}[plain]
\begin{tikzpicture}[remember picture,overlay]
  \fill[Cream] (current page.south west) rectangle (current page.north east);
  \fill[Navy] (current page.south west) rectangle ([yshift=1.15cm]current page.south east);
  \fill[SoftBlue] ([xshift=1.0cm,yshift=-1.0cm]current page.north west) circle (0.85cm);
\end{tikzpicture}

\vspace{0.9cm}
\begin{center}
{\color{Navy}\fontsize{24}{28}\selectfont\bfseries Terima Kasih}\\[0.2cm]
{\color{TextDark}\large NechCode -- Ide Bisnis dan Analisa Kelayakan}\\[0.35cm]
\pill{Kelompok 12}
\end{center}

\vspace{0.4cm}
\begin{columns}[T,totalwidth=\textwidth]
\column{0.48\textwidth}
\begin{block}{Anggota}
\small
Louis Cristian S. L -- 5002221004\\
Hafidz Mulia -- 5002221022\\
Nabeel Farvez F. -- 5002221139
\end{block}

\column{0.48\textwidth}
\begin{block}{Anggota}
\small
M. Rajwa Dzaki D. -- 5002221143\\
Stanadyne Tolanda -- 5046231038
\end{block}
\end{columns}

\vfill

\end{frame}

\end{document}